\documentclass[12pt,letterpaper]{article}

% ===================== PAQUETES =====================
\usepackage[utf8]{inputenc}
\usepackage[spanish]{babel}
\usepackage[T1]{fontenc}
\usepackage{times}
\usepackage[margin=2.5cm]{geometry}
\usepackage{setspace}
\usepackage{titlesec}
\usepackage{tocloft}
\usepackage{fancyhdr}
\usepackage{hyperref}
\usepackage{tcolorbox}
\usepackage{xcolor}
\usepackage[numbers,sort]{natbib}

% ===================== CONFIGURACIÓN GENERAL =====================
\onehalfspacing
\hypersetup{
  colorlinks=true, linkcolor=black, urlcolor=blue, citecolor=black,
  pdftitle={Precisión real versus percepción estudiantil: análisis comparativo de las estimaciones de tiempo de viaje de Google Maps en usuarios frecuentes de transporte urbano en la UPTC Tunja},
  pdfauthor={Paula Andrea Betancourt Matamoros, Brayan Stiven Buitrago Pulido, Maira Suleima Rodríguez González},
  pdfsubject={Metodología de la Investigación y Diseño Experimental},
}


% Numeración de secciones tipo "1.", "1.1", etc.
\renewcommand{\thesection}{\arabic{section}}
\renewcommand{\thesubsection}{\thesection.\arabic{subsection}}

\titleformat{\section}{\centering\bfseries\large}{\thesection.}{0.5em}{}
\titleformat{\subsection}{\bfseries\large}{\thesubsection}{0.5em}{}

% Nombres en español
\renewcommand{\contentsname}{Tabla de contenido}
\renewcommand{\listtablename}{Índice de tablas}
\renewcommand{\listfigurename}{Índice de figuras}
\renewcommand{\tablename}{Tabla}
\renewcommand{\figurename}{Figura}

% Caja para el estado "EN CONSTRUCCIÓN"
\newtcbox{\enconstruccion}{on line, boxrule=0.4pt, colframe=orange!70!black,
  colback=orange!5, arc=6pt, boxsep=1pt, left=4pt, right=4pt, top=1pt, bottom=1pt,
  fontupper=\scriptsize\bfseries\color{orange!70!black}}

\newcommand{\estado}[1]{\begin{center}\enconstruccion{#1}\end{center}}
\newcommand{\notaseccion}[1]{\begin{center}\itshape\footnotesize #1\end{center}}

% ===================== DATOS DEL DOCUMENTO =====================
\newcommand{\tituloProvisional}{Precisión real versus percepción estudiantil: análisis comparativo de las estimaciones de tiempo de viaje de Google Maps en usuarios frecuentes de transporte urbano en la UPTC Tunja}

\begin{document}

% ===================== PORTADA =====================
\begin{titlepage}
\centering
\vspace*{1cm}
{\large UNIVERSIDAD PEDAGÓGICA Y TECNOLÓGICA DE COLOMBIA\par}
{\large FACULTAD DE INGENIERÍA\par}
{\large INGENIERÍA DE SISTEMAS Y COMPUTACIÓN\par}

\vspace{3cm}
{\LARGE\bfseries \tituloProvisional \par}

\vspace{0.5cm}
{\itshape Título provisional --- se ajusta al formular la propuesta (Fase 3)\par}

\vspace{3cm}
{\large AUTORES\par}
\vspace{0.5cm}
{\large Paula Andrea Betancourt Matamoros\par}
{\large Brayan Stiven Buitrago Pulido\par}
{\large Maira Suleima Rodríguez González\par}

\vfill
{\itshape Documento elaborado en el marco del curso Metodología de la Investigación\par}
\vspace{0.5cm}
{Tunja, Colombia · 2026\par}
\end{titlepage}

% ===================== TABLA DE CONTENIDO =====================
\tableofcontents
\newpage

% ===================== ÍNDICE DE TABLAS =====================
\listoftables
\notaseccion{Aún no hay tablas incorporadas al cuerpo del documento. Se listan aquí a medida
que cada fase incorpore tablas propias (matrices, resultados, operacionalización de
variables, etc.).}
\newpage

% ===================== ÍNDICE DE FIGURAS =====================
\listoffigures
\notaseccion{Aún no hay figuras incorporadas al cuerpo del documento. Se listan aquí a medida
que cada fase incorpore figuras propias (diagramas, gráficos, esquemas, etc.).}
\newpage

% ===================== SECCIONES =====================
\section{Planteamiento del problema}

\subsection{Antecedentes}

Google Maps se ha convertido en una herramienta cotidiana para que los estudiantes
universitarios planeen sus desplazamientos. Su función principal es indicar dónde se
encuentra un lugar y calcular cuánto tiempo tomará llegar a él, combinando para ello
distancia, velocidad promedio y datos de tráfico \cite{salvatierra2023}.

Estimar ese tiempo con precisión es más complejo que calcular una distancia en
kilómetros, porque depende de condiciones viales que cambian constantemente. El
algoritmo de Google Maps se basa principalmente en patrones históricos de tráfico, lo
que limita su capacidad de anticipar sucesos puntuales —accidentes, obras,
trancones— que aparecen sin previo aviso \cite{thanuja2026}.

Este algoritmo se calibra, además, con datos recolectados principalmente en ciudades
grandes, donde el volumen de usuarios conectados y el tipo de infraestructura vial
generan patrones de tráfico distintos a los de una ciudad intermedia como Tunja; el
desarrollo de soluciones de movilidad inteligente suele orientarse a contextos de
mayor escala urbana, y su desempeño en ciudades intermedias no está documentado con
el mismo nivel de detalle \cite{pulido2024}.

Tunja presenta, a su vez, condiciones viales que agravan esta diferencia: calles
angostas en el centro histórico, desvíos por obras que no se reportan con
anticipación y una operación del transporte público colectivo urbano cuya
regularidad y organización han sido cuestionadas en estudios previos
\cite{munoz2019}. A esto se suman prácticas propias de la operación diaria de las
busetas, como paradas no programadas para recoger o dejar pasajeros, rutas alteradas
por decisión del conductor y tiempos de espera irregulares entre unidades, factores
que Google Maps no incorpora en tiempo real.

Como consecuencia de estas limitaciones, el tiempo real que un estudiante invierte en
su trayecto entre su vivienda y la universidad suele diferir del tiempo que Google
Maps había calculado previamente. Para quienes usan el transporte urbano colectivo al
menos tres veces por semana, esta discrepancia entre lo estimado y lo vivido no es un
hecho aislado, sino una situación que se repite semana tras semana.

El primer efecto de esta discrepancia es la llegada tarde a clase, incluso cuando el
estudiante organizó su salida siguiendo de forma estricta el tiempo indicado por la
aplicación. La repetición de este error genera, además, estrés o ansiedad antes de
ingresar a clase, una respuesta que se explica por la exigencia del entorno
universitario frente a situaciones que el estudiante percibe como fuera de su
control.

Cuando el margen de error se acumula, las consecuencias dejan de ser solo
emocionales: el estudiante puede perder clases, recibir llamados de atención o ver
afectada una evaluación, lo que impacta directamente su rendimiento académico.

La reiteración de estas fallas, semana tras semana, hace que el estudiante deje de
confiar específicamente en la precisión del tiempo de llegada calculado por Google
Maps —no en la aplicación como tal ni en su función de ubicación— y deje de
considerarla una herramienta confiable para planear con exactitud la hora de salida
hacia la universidad \cite{peer2025}.

Ante esta pérdida de confianza, varios estudiantes adoptan un margen de seguridad:
salen de su casa antes de lo indicado "por si acaso", o cambian de ruta o de medio de
transporte de forma improvisada para compensar el margen de error que la aplicación
no corrige. Este comportamiento compensatorio surge, en parte, del desconocimiento
que tiene el algoritmo sobre los horarios pico reales de Tunja y sobre las rutas
alternas que solo identifica quien usa el transporte urbano de manera cotidiana
\cite{kapatsila2025}.

En conjunto, estos hallazgos evidencian una brecha no documentada entre el desempeño
de Google Maps para estimar tiempos de viaje en ciudades grandes y su desempeño en
una ciudad intermedia como Tunja. Esta brecha afecta de manera concreta la vida
académica de los estudiantes de la UPTC, sede central, que usan el transporte
público urbano al menos tres veces por semana para llegar a clase, y constituye el
problema que la presente investigación propone estudiar.

\subsection{Descripción de la situación problemática}
\estado{EN CONSTRUCCIÓN}
\notaseccion{Se redacta cuando se libere el tema correspondiente en el Curso.}

\subsection{Formulación del problema}
\estado{EN CONSTRUCCIÓN}
\notaseccion{Se redacta cuando se libere el tema correspondiente en el Curso.}

\subsection{Consecuencias de no investigar}
\estado{EN CONSTRUCCIÓN}
\notaseccion{Se redacta cuando se libere el tema correspondiente en el Curso.}

\subsection{Delimitación}
\estado{EN CONSTRUCCIÓN}
\notaseccion{Se redacta cuando se libere el tema correspondiente en el Curso.}
\newpage
\section{Justificación}
\estado{EN CONSTRUCCIÓN}
\notaseccion{Se redacta en la Fase 3 del Curso (Formulación de la propuesta).}
\newpage
\section{Objetivos}
\estado{EN CONSTRUCCIÓN}
\notaseccion{Se redacta en la Fase 3 del Curso.}
\newpage
\section{Alcance de la investigación}
\estado{EN CONSTRUCCIÓN}
\notaseccion{Sección incorporada según Hernández-Sampieri (tipo de estudio: exploratorio,
descriptivo, correlacional o explicativo) --- no aparece como título propio en la NTC 1486.
Se redacta en la Fase 3.}
\newpage
\section{Formulación de hipótesis}
\estado{EN CONSTRUCCIÓN}
\notaseccion{Sección incorporada según Hernández-Sampieri, Tamayo y Tamayo y Bernal ---
no aparece como título propio en la NTC 1486. Se redacta en la Fase 3.}
\newpage
\section{Variables}
\estado{EN CONSTRUCCIÓN}
\notaseccion{Definición nominal/conceptual en la Fase 3; operacionalización completa en la Fase 5.}
\newpage
\section{Marco referencial}
\estado{EN CONSTRUCCIÓN}
\notaseccion{Se redacta en la Fase 4 del Curso (Construcción del marco referencial).}
\newpage
\section{Diseño metodológico}
\estado{EN CONSTRUCCIÓN}
\notaseccion{Se redacta en la Fase 5 del Curso.}
\newpage
\section{Resultados y discusión}
\estado{EN CONSTRUCCIÓN}
\notaseccion{Se redacta en la Fase 6 del Curso.}
\newpage
\section{Conclusiones y recomendaciones}
\estado{EN CONSTRUCCIÓN}
\notaseccion{Se redacta en la Fase 7 del Curso.}
\newpage

% ===================== REFERENCIAS =====================
\section*{}
\addcontentsline{toc}{section}{Referencias}
\bibliographystyle{ieeetr}
\bibliography{referencias}
\notaseccion{Esta lista crece con cada sección que se redacte; la numeración se reajusta
para mantener el orden de primera aparición en todo el documento.}
\newpage

% ===================== ANEXOS =====================
\section*{Anexos}
\addcontentsline{toc}{section}{Anexos}
\estado{EN CONSTRUCCIÓN}
\notaseccion{Reúne los instrumentos usados en cada fase (matriz de correlación, ficha de
causas filtradas, diagrama de Ishikawa, y los que se sumen más adelante).}

\end{document}